%%=============================================================================
%% Inleiding
%%=============================================================================

\chapter{\IfLanguageName{dutch}{Inleiding}{Introduction}}%
\label{ch:inleiding}

Netwerkbeheer vormt een kritieke component binnen moderne IT-infrastructuren. Het accuraat bijhouden van netwerkinventarisaties en -documentatie is essentieel voor efficiënt beheer, probleemoplossing en beveiliging. Echter, in veel organisaties raken deze inventarisaties snel verouderd omdat wijzigingen in de infrastructuur niet consistent worden bijgehouden. Dit probleem manifesteert zich vooral in dynamische IT-omgevingen waar systemen en apparaten voortdurend veranderen, waardoor handmatige documentatieprocessen onvoldoende zijn om de actualiteit te garanderen.\autocite{NetBoxPlanning}

Dit onderzoek wordt uitgevoerd binnen de context van Citymesh, een Belgische telecomoperator en technologiebedrijf dat gespecialiseerd is in grootschalige draadloze netwerken en connectiviteitsoplossingen voor de zakelijke markt.\autocite{CitymeshAbout} Citymesh focust op zowel permanente als tijdelijke connectiviteit op basis van WiFi/WLAN, 0G/IoT, 4G en 5G en treedt op als vierde mobiele operator in België.\autocite{CitymeshAbout,Mice2022} Deze heterogene en sterk gedistribueerde draadloze infrastructuren veranderen frequent door nieuwe sites, tijdelijke events en uitbreidingen, waardoor het bijhouden van accurate netwerkdocumentatie een bijzondere uitdaging vormt. Citymesh heeft specifiek gevraagd om een evaluatie en implementatie van NetBox Discovery als oplossing voor deze uitdaging. Binnen die bredere context wordt de praktische uitwerking in deze bachelorproef bewust afgebakend tot een casestudy op concrete sites en op documentatie-elementen die in NetBox rechtstreeks toetsbaar waren, met name inventarisdata en kabel-/topologierelaties.

De doelgroep van dit onderzoek bestaat uit netwerkbeheerders en infrastructuurbeheerders binnen organisaties die te maken hebben met verouderde of inconsistente netwerkdocumentatie. Specifiek richt dit onderzoek zich op professionals binnen Citymesh die verantwoordelijk zijn voor het onderhoud van netwerkinventarisaties in complexe draadloze omgevingen en die behoefte hebben aan geautomatiseerde oplossingen om een betrouwbare ``single source of truth'' te realiseren.

\section{\IfLanguageName{dutch}{Probleemstelling}{Problem Statement}}%
\label{sec:probleemstelling}

Het centrale probleem dat dit onderzoek adresseert, is de snel verouderende netwerkdocumentatie binnen organisaties. Wanneer wijzigingen in de infrastructuur niet consistent worden bijgehouden, ontstaan onnauwkeurige of tegenstrijdige netwerkdocumentatie. Dit leidt tot het ontbreken van een betrouwbare single source of truth (SSOT), wat op zijn beurt risico's vormt voor efficiënt beheer, probleemoplossing en beveiliging.\autocite{NetBoxOSS} Het probleem wordt versterkt in dynamische IT-omgevingen zoals die van Citymesh, waar draadloze technologieën zoals grootschalige WiFi/WLAN, private 4G/5G-netwerken en IoT- of 0G-sensornetwerken continu evolueren en waar apparaten en sites frequent toegevoegd, verhuisd of verwijderd worden.\autocite{CitymeshAbout,CitymeshNetworkService}

\section{\IfLanguageName{dutch}{Onderzoeksvraag}{Research question}}%
\label{sec:onderzoeksvraag}

\subsection{Centrale onderzoeksvraag}

De hoofdonderzoeksvraag van dit onderzoek luidt:

\textit{Hoe kan automatische netwerkdetectie bijdragen aan het up-to-date houden van netwerkdocumentatie en het realiseren van een betrouwbare `single source of truth' binnen de dynamische, draadloze IT-infrastructuren van Citymesh?}

\subsection{Deelonderzoeksvragen}

Om de centrale onderzoeksvraag te beantwoorden, worden de volgende deelonderzoeksvragen geformuleerd:

\textbf{Probleemdomeinvragen:}
\begin{enumerate}
  \item Hoe manifesteert het probleem van verouderde netwerkdocumentatie zich specifiek binnen de draadloze (WiFi/WLAN, private 4G/5G, IoT/0G) infrastructuur van Citymesh?
  \item Wat zijn de oorzaken van verouderde netwerkdocumentatie bij Citymesh, en welke factoren dragen hieraan bij?
  \item Wat is de huidige staat van de netwerkdocumentatie bij Citymesh, en wat zijn de gevolgen van verouderde of inconsistente documentatie voor het netwerkbeheer?
  \item Hoe wordt netwerkdocumentatie momenteel bijgehouden bij Citymesh, en welke processen en tools worden daarbij gebruikt?
\end{enumerate}

\textbf{Oplossingsdomeinvragen:}
\begin{enumerate}
  \setcounter{enumi}{4}
  \item Welke bestaande technieken en tools voor automatische netwerkdetectie bestaan er, en in welke mate ondersteunen zij grootschalige draadloze omgevingen zoals die van Citymesh?
  \item Hoe presteren geselecteerde automatische netwerkdetectietools (met focus op NetBox Discovery) qua nauwkeurigheid, dekking, integratiemogelijkheden en beheerlast in de context van Citymesh?
  \item In welke mate verbetert NetBox Discovery de actualiteit, foutgraad en onderhoudstijd van de netwerkdocumentatie ten opzichte van de huidige werkwijze?
  \item Welke organisatorische processen, best practices en configuratie-aanpassingen zijn nodig om een duurzame SSOT voor netwerkdocumentatie te waarborgen met NetBox Discovery?
\end{enumerate}

\section{\IfLanguageName{dutch}{Onderzoeksdoelstelling}{Research objective}}%
\label{sec:onderzoeksdoelstelling}

Het concrete eindresultaat van dit onderzoek bestaat uit:
\begin{itemize}
  \item Een evaluatie van NetBox Discovery, inclusief een vergelijking met alternatieve automatische netwerkdetectietools en hun effectiviteit in het bijhouden van netwerkdocumentatie
  \item Een praktische implementatie en configuratie van NetBox Discovery binnen de draadloze infrastructuur van Citymesh
  \item Een rapport met concrete voor/na-observaties en gerichte metingen over de verbetering in actualiteit en betrouwbaarheid van netwerkdocumentatie na implementatie van NetBox Discovery
  \item Concrete aanbevelingen voor Citymesh over de implementatie, configuratie en organisatorische inbedding van NetBox Discovery
\end{itemize}

De bachelorproef kan als succes beschouwd worden wanneer een werkende implementatie gerealiseerd is die aantoonbaar bijdraagt aan het verbeteren van de actualiteit van netwerkdocumentatie, en wanneer concrete observaties en metingen beschikbaar zijn die de meerwaarde van automatische detectie binnen de afgebakende casus onderbouwen.

\section{\IfLanguageName{dutch}{Opzet van deze bachelorproef}{Structure of this bachelor thesis}}%
\label{sec:opzet-bachelorproef}

De rest van deze bachelorproef is als volgt opgebouwd:

In Hoofdstuk~\ref{ch:stand-van-zaken} wordt een overzicht gegeven van de stand van zaken binnen het onderzoeksdomein, op basis van een literatuurstudie. Hierbij wordt ingegaan op het concept van single source of truth in netwerkbeheer, automatische netwerkdetectie en bestaand onderzoek naar automatische documentatie.

In Hoofdstuk~\ref{ch:methodologie} wordt de methodologie toegelicht: de zes onderzoeksfasen, gebruikte tools en technologieën, en hoe de onderzoeksvragen worden beantwoord.

Daarna volgen hoofdstukken die de uitwerking van die fasen beschrijven. In Hoofdstuk~\ref{ch:fase1} staat de requirementsanalyse bij Citymesh. Hoofdstuk~\ref{ch:fase2} behandelt de toolselectie en vergelijkende evaluatie. Hoofdstuk~\ref{ch:fase3} beschrijft de baselinemetingen en de voorbereiding van de implementatie. Hoofdstuk~\ref{ch:fase4} gaat over de implementatie en configuratie van NetBox Discovery (Orb agent, Diode). Hoofdstuk~\ref{ch:fase5} behandelt de evaluatie en metingen. Hoofdstuk~\ref{ch:fase6} bevat analyse, aanbevelingen en reflectie.

In Hoofdstuk~\ref{ch:conclusie}, tenslotte, wordt de conclusie gegeven en een antwoord geformuleerd op de onderzoeksvragen. Daarbij wordt ook een aanzet gegeven voor toekomstig onderzoek binnen dit domein.
